% Generated from validated Article Draft Result. Do not hand-edit; revise the structured draft instead.
\section{2023年総括――「モデル」から「システム／生態系」へ}
\noindent\textbf{2023年を一つの年として読むと、生成AIの設計対象はモデルweightsだけに留まらなくなった。インタラクション、公開モデル群、tool/action、runtime、retrieval/post-training、生成メディア、evaluation/controlが別々の層として立ち上がり、相互の制約を受けるシステムとして見えるようになった。}

年初には「より強い基盤モデル」が議論の中心に見えやすかったが、年を通じて利用形態と実装層が増えた。年末を同じ尺度で見ると、モデルの出力品質だけでなく、どの入力を扱い、どの外部機能へ接続し、どの資源で運用し、どう評価・制御するかまでが設計問題として分化している。


この変化は、2023年の出来事を「モデル発表」「論文」「画像生成」といった媒体別に並べるだけでは捉えにくい。annual scaleでは、同じモデル系列でもinterface、deployment、runtime、controlの異なる層にまたがるため、生成AIを多層system/ecosystemとして再分類する方が技術的な因果を説明しやすい。


open weightsの拡大はruntime・fine-tuning効率への需要を強め、tool interfaceの整備はreliabilityとgovernanceの境界を明確にし、長文脈やmultimodal化は評価方法を複雑にした。各層は独立に発展したのではなく、別の層へ新しい制約を渡しながら設計空間を広げた。


一方、2023年末でも未解決の境界は多い。公称context長と実効利用、vendor benchmarkと独立評価、open weightsと再現可能なopen training、tool selectionと安全な実行、生成品質と権利・安全性、capability評価と実効的なgovernanceは同一視できない。


したがって2023年の年次的な転換は「ある一つのモデルが勝った」ことではなく、生成AIがモデル単体の競争から、多層の技術スタックと生態系を設計する競争へ移ったことにある。この見取り図が、翌年以降のon-device、agent、serving、multimodal、controlの加速を理解する基準になる。
